\section*{Capítulo \ref{ch:b1:electrophilic-additions} --- Alquenos y alquinos: adiciones electrófilas}

\begin{solution}{exo:b1:electrophilic-additions:1}
2-Cloropropano; 2-bromo-2-metilpropano; 1-bromo-1-metilciclohexano;
2-clorobut-2-eno.
\end{solution}

\begin{solution}{exo:b1:electrophilic-additions:2}
1,2-Dibromoetano; \textit{trans}-1,2-dibromociclohexano;
1,2-dibromopropano.
\end{solution}

\begin{solution}{exo:b1:electrophilic-additions:3}
Propan-2-ol; 2-metilpropan-2-ol; 1-metilciclohexan-1-ol (Markovnikov en
cada caso).
\end{solution}

\begin{solution}{exo:b1:electrophilic-additions:4}
Propino (\omterm{def:b1:electrophilic-additions:acetylide}{alquino terminal}): \ce{CH3C#CH + NH2- -> CH3C#C- + NH3}. Etanol:
\ce{CH3CH2OH + NH2- -> CH3CH2O- + NH3}. El but-2-ino y el propeno no
tienen ningún hidrógeno lo bastante ácido.
\end{solution}

\begin{solution}{exo:b1:electrophilic-additions:5}
El par $\pi$ toma \ce{H+} sobre el carbono \ce{CH2}: catión terciario
\ce{(CH3)3C+}; un \omterm{def:b1:lewis-resonance-vsepr:lewis}{par no enlazante} del agua se une a él:
\ce{(CH3)3COH2+}; una molécula de agua toma un protón: \ce{(CH3)3COH}. En
ácido concentrado caliente, el agua escasea y el alqueno volátil escapa:
avanza la reacción inversa, la \omterm{def:b1:alcohol-activation:dehydration}{deshidratación}.
\end{solution}

\begin{solution}{exo:b1:electrophilic-additions:6}
La protonación del 2-metilpropeno da un catión terciario, y la del
propeno, uno secundario. La protonación, \omterm{def:b1:elementary-steps:rds}{etapa determinante de la
velocidad}, tiene un \omterm{def:b1:elementary-steps:energy-profile}{estado de transición} que se parece al catión
(Hammond): al catión terciario, más estable, se llega por una barrera más
baja.
\end{solution}

\begin{solution}{exo:b1:electrophilic-additions:7}
El ion bromonio forma un puente sobre una cara del anillo; el bromuro lo
abre por la otra cara: los dos bromos terminan en \textit{trans}. El
dibromuro \textit{trans} es \omterm{def:b1:stereochemistry-in-depth:stereocentre}{quiral} (sin plano de simetría), pero el
bromuro ataca por igual los dos carbonos: los dos \omterm{def:b1:stereochemistry-in-depth:enantiomers}{enantiómeros} se forman
en cantidades iguales, y el producto es racémico, ópticamente inactivo.
\end{solution}

\begin{solution}{exo:b1:electrophilic-additions:8}
La protonación del C1 da el catión secundario en el C2, contiguo al C3
terciario que lleva un hidrógeno; ese hidrógeno migra con su par al C2
(una migración 1,2 de hidruro) y deja un catión terciario en el C3, que
captura el \ce{Cl-}: 2-cloro-2-metilbutano.
\end{solution}

\begin{solution}{exo:b1:electrophilic-additions:9}
Propino + amida de sodio: el ion propinuro; con 1-bromopropano (SN2):
\ce{CH3C#C-CH2CH2CH3}, hex-2-ino. El propinuro con el benzaldehído y,
después, ácido diluido: \ce{C6H5CH(OH)C#CCH3}, 1-fenilbut-2-in-1-ol.
\end{solution}

\begin{solution}{exo:b1:electrophilic-additions:10}
\cip{E}-but-2-eno: ion bromonio sobre una cara, bromuro desde la otra en
el C2 o en el C3; ambos ataques dan el mismo compuesto, de descriptores
\cip{2R,3S}: el \omterm{def:b1:stereochemistry-in-depth:enantiomers}{compuesto meso}, un solo estereoisómero.
\cip{Z}-but-2-eno: el ataque al C2 da \cip{2S,3S} (o su imagen especular
desde la otra cara), y el ataque al C3, \cip{2R,3R}; en cantidades
iguales: dos \omterm{def:b1:stereochemistry-in-depth:isomers}{estereoisómeros}, una \omterm{def:b1:stereochemistry-in-depth:enantiomers}{mezcla racémica}.
\end{solution}

\begin{solution}{exo:b1:electrophilic-additions:11}
En el ion bromonio asimétrico, el carbono más sustituido soporta una
mayor parte de la carga positiva (su enlace C--Br es más largo y más
débil, como un catión casi terciario o secundario). El agua, en gran
exceso, ataca ese carbono desde la cara opuesta al puente:
1-bromopropan-2-ol.
\end{solution}

\begin{solution}{exo:b1:electrophilic-additions:12}
$D = 1$: un alqueno. El 2-metilbut-1-eno y el 2-metilbut-2-eno dan ambos
el catión terciario y, por tanto, 2-bromo-2-metilbutano y
2-metilbutan-2-ol. Sus espectros de RMN de protón difieren: dos protones
vinílicos (\ce{=CH2}, singuletes) en el primero, uno (\ce{=CH-}, un
cuadruplete) en el segundo.
\end{solution}

\begin{solution}{pb:b1:electrophilic-additions:1}
\textbf{1.} En cada carbono, \ce{CH3} tiene prioridad sobre H: el \cip{Z}
tiene los dos metilos del mismo lado, el \cip{E}, en lados opuestos.
\textbf{2.} \omterm{def:b1:stereochemistry-in-depth:isomers}{Estereoisómeros} que no son imágenes especulares:
\omterm{def:b1:stereochemistry-in-depth:enantiomers}{diastereoisómeros}.
\textbf{3.} La rotación en torno al C=C rompería el enlace $\pi$, lo que
cuesta mucha más energía de la que aportan los choques a temperatura
ambiente.
\textbf{4.} El \cip{E}: sus grupos metilo no se estorban.
\textbf{5.} \ce{C4H8 + Br2 -> C4H8Br2}.
\textbf{6.} Dos \omterm{def:b1:stereochemistry-in-depth:stereocentre}{centros estereogénicos} equivalentes: como máximo cuatro,
en realidad tres (\cip{2R,3R}, \cip{2S,3S} y el meso \cip{2R,3S}).
\textbf{7.} El par $\pi$ ataca un átomo de Br del \ce{Br2}, el par Br--Br
sale como \ce{Br-}, y un \omterm{def:b1:lewis-resonance-vsepr:lewis}{par no enlazante} del bromo atacado se une al
segundo carbono: un ion bromonio de tres miembros.
\textbf{8.} En el ion con puente todos los átomos tienen octeto; un
\omterm{def:b1:electronic-effects:intermediates}{carbocatión} secundario abierto tendría un carbono con seis electrones.
\textbf{9.} Un \omterm{def:b1:lewis-resonance-vsepr:lewis}{par no enlazante} del \ce{Br-} se une a un carbono del
anillo; el enlace C--Br del puente se rompe, y su par pasa al bromo del
puente.
\textbf{10.} El puente ocupa una cara; como en la SN2, el \omterm{def:b1:electronic-effects:nucleophile}{nucleófilo} entra
por el lado opuesto al enlace que se rompe.
\textbf{11.} Una \omterm{def:b1:electrophilic-additions:halonium}{adición anti}.
\textbf{12.} El agua competiría con el bromuro por el ion bromonio y daría
una bromohidrina.
\textbf{13.} \cip{2R,3S} (o \cip{2S,3R}, el mismo compuesto).
\textbf{14.} El mismo compuesto.
\textbf{15.} En una \omterm{def:b1:stereochemistry-in-depth:isomers}{conformación} con los dos bromos eclipsados, un plano
de simetría pasa entre el C2 y el C3; los descriptores son opuestos.
\textbf{16.} \cip{2R,3R} y \cip{2S,3S}.
\textbf{17.} En cantidades iguales: una \omterm{def:b1:stereochemistry-in-depth:enantiomers}{mezcla racémica}.
\textbf{18.} Sí: los dos alquenos \omterm{def:b1:stereochemistry-in-depth:enantiomers}{diastereoisómeros} dan productos
distintos (meso frente a racémico).
\textbf{19.} Ninguna: los dos \omterm{def:b1:stereochemistry-in-depth:enantiomers}{enantiómeros} se compensan.
\textbf{20.} $5.60/56.0 = \qty{0.100}{mol}$ de cada uno.
\textbf{21.} $0.100 \times 0.85 \times 215.8 = \qty{18.3}{g}$ de cada uno.
\textbf{22.} \textbf{$[\alpha] = 0^\circ$ (\omterm{def:b1:stereochemistry-in-depth:enantiomers}{compuesto meso}), \qty{18.3}{g}}.
\end{solution}
